\begin{textsample}{POS Dim 8 – human – Score 2.00 – mg/mg\_ef\_er\_7\_p5.txt}  \label{ex:f8_pos_018}
Nesse componente curricular não se trata de ensinar uma ou diversas religiões, mas oportunizar que os educandos reflitam sobre suas experiências, vivências, sobre as culturas e tradições religiosas, aprofundando e formulando sentidos para a construção de seus projetos de vida com qualidade, capacidade de respeito, em vista de uma cultura solidária e de paz.
Destaca -se que os critérios de organização das habilidades no CRMG, a partir da BNCC, com a explicitação dos objetos de conhecimento, aos quais se relacionam, e do agrupamento desses objetos em unidades temáticas, construídos com a participação e a contribuição de centenas de professores, expressam um arranjo possível, \textbf{dentre} outros.
Portanto, os agrupamentos propostos são uma referência básica, importante, e não devem ser tomados como modelo fixo e obrigatório, o que possibilita a adaptação, o aperfeiçoamento constante, a partir dos princípios, competências gerais e específicas da Base, e o atendimento à \textbf{realidade} dos educandos e de \textbf{realidades} locais e \textbf{regionais}.

% matched lemmas: dentrar, realidade, regional
\end{textsample}
